% ======================================================================
%  PART 1 -- Foundations: Stationarity, Autocovariance, Ergodicity
% ======================================================================
\section{Foundations: Stationarity, Autocovariance \& Ergodicity}
\label{sec:foundations}

A \emph{time series} $\{X_t\}$ is both the data recorded in time order and the
stochastic process $\{X_t : t\in T\}$ on a probability space
$(\Omega,\mathcal F,\Prob)$ it is a realisation of. The index set $T$ is a set
of time points (here $T=\Ints$ unless stated otherwise; we take the sampling
step $\Delta=1$). Time series differ from classical statistics because the
observations are \textbf{dependent}: ``there are many more ways to be dependent
than to be independent''. Everything in this course is built on describing that
dependence through second-order structure.

% ----------------------------------------------------------------------
\subsection{Definitions}

\begin{definition}{(Weak / second-order) stationarity}{wkstat}
$\{X_t\}$ is \textbf{second-order} (equivalently \textbf{weakly} or
\textbf{covariance}) \textbf{stationary} if all first- and second-order joint
moments exist, are finite, and are invariant to time shifts. Concretely, for
all $t,s\in T$ and all admissible shifts $\tau$:
\[
\text{(i) } \E[X_t]=\mu;\qquad
\text{(ii) } \Var(X_t)=\sigma^2<\infty;\qquad
\text{(iii) } \E[X_tX_{t+\tau}]=\E[X_sX_{s+\tau}].
\]
Condition (iii) forces $\E[X_tX_{t+\tau}]$ to be a function of the lag $\tau$
only.
\end{definition}

\begin{definition}{Strong / strict stationarity}{ststat}
$\{X_t\}$ is \textbf{strictly} (completely) \textbf{stationary} if for every
$n\ge1$, every $t_1,\dots,t_n$ and every shift $\tau$, the joint distribution
of $(X_{t_1},\dots,X_{t_n})$ equals that of
$(X_{t_1+\tau},\dots,X_{t_n+\tau})$.
\end{definition}

\begin{definition}{Gaussian process}{gauss:gp}
$\{X_t\}$ is a \textbf{Gaussian process} if for all $n$ and $t_1,\dots,t_n$ the
vector $(X_{t_1},\dots,X_{t_n})$ is multivariate normal. Such a process is
completely characterised by its (finite) first two moments; hence for Gaussian
processes \textbf{weak $\Leftrightarrow$ strict stationarity}.
\end{definition}

\begin{definition}{Autocovariance sequence (ACVS)}{acvs}
For a discrete-time second-order stationary $\{X_t\}$, the \textbf{ACVS} is
\[
\acvs_\tau=\Cov(X_0,X_\tau)=\Cov(X_t,X_{t+\tau}),\qquad \tau\in\Ints .
\]
Immediate properties: $\acvs_0=\Var(X_t)\ge0$, symmetry $\acvs_{-\tau}=\acvs_\tau$,
and (Proposition~\ref{prop:psd}) positive semi-definiteness.
\end{definition}

\begin{definition}{Autocorrelation sequence (ACF)}{acf}
The \textbf{ACF} is $\acf_\tau=\Corr(X_t,X_{t+\tau})=\acvs_\tau/\acvs_0$, with
$\acf_0=1$ and $\abs{\acf_\tau}\le1$.
\end{definition}

\begin{definition}{Sample autocovariance estimators}{sampleacvs}
Given $X_1,\dots,X_n$ with sample mean $\bar X$, two estimators are used for
$\abs{\tau}\le n-1$ (both $0$ otherwise):
\[
\tilde\acvs_\tau=\frac{1}{n-\abs\tau}\sum_{t=1}^{n-\abs\tau}(X_t-\bar X)(X_{t+\abs\tau}-\bar X),
\qquad
\hat\acvs_\tau=\frac{1}{n}\sum_{t=1}^{n-\abs\tau}(X_t-\bar X)(X_{t+\abs\tau}-\bar X).
\]
With known mean, $\tilde\acvs_\tau$ is \emph{unbiased} but need not be
positive semi-definite; $\hat\acvs_\tau$ is \emph{biased}
($\E[\hat\acvs_\tau]=\tfrac{n-\abs\tau}{n}\acvs_\tau$) but \emph{is} positive
semi-definite, which is why it is the default. The plug-in ACF estimator is
$\hat\acf_\tau=\hat\acvs_\tau/\hat\acvs_0$.
\end{definition}

\begin{definition}{Mean ergodicity}{ergodic}
$\{X_t\}$ is \textbf{mean ergodic} if its first two moments are finite and
$\bar X\toprob \E[X_t]$ as $n\to\infty$ (``temporal averages
$=$ population averages''). Ergodicity and stationarity are \emph{not}
equivalent.
\end{definition}

% ----------------------------------------------------------------------
\subsection{Theorems \& Propositions}

\begin{theorem}{Kolmogorov existence theorem}{kolmo}
A family of finite-dimensional distribution functions $\{F_{\mathbf t}\}$ are
the f.d.d.'s of a stochastic process \emph{iff} they are
\textbf{consistent}: for every $n$, every $\mathbf t$ and every $1\le i\le n$,
\[
\lim_{x_i\to\infty}F_{\mathbf t}(\mathbf x)=F_{\mathbf t(i)}(\mathbf x(i)),
\]
where $\mathbf t(i),\mathbf x(i)$ delete the $i$-th coordinate.\\[2pt]
\textit{Why:} it guarantees an infinite-index stochastic process actually
exists given only its finite-dimensional ``marginals'', justifying modelling
$\{X_t\}_{t\in\Ints}$ even though we observe finitely many points.
\end{theorem}

\begin{proposition}{The ACVS is positive semi-definite}{psd}
For any $n$, any $t_1,\dots,t_n\in\Ints$ and any $a_1,\dots,a_n\in\Reals$,
\[
\sum_{j=1}^n\sum_{k=1}^n a_ja_k\,\acvs_{j-k}\ge0 .
\]
\textit{Proof idea:} set $W=\sum_j a_jX_j$; then
$0\le\Var(W)=\sum_{j,k}a_ja_k\Cov(X_j,X_k)=\sum_{j,k}a_ja_k\acvs_{j-k}$.\\
\textit{Why:} p.s.d.\ is the defining algebraic property of a valid ACVS; a
candidate sequence that is \emph{not} p.s.d.\ \textbf{cannot} be an
autocovariance (used constantly to reject ``fake'' ACVS, see
Exercise in Part~\ref{sec:spectral}).
\end{proposition}

\begin{proposition}{Variance of the sample mean \& consistency}{xbarvar}
If $\sum_{\tau}\abs{\acvs_\tau}<\infty$ then $\bar X$ is unbiased and
\[
\Var(\bar X)=\frac1{n^2}\sum_{i,j=1}^n\acvs_{j-i}
=\frac1{n}\sum_{\tau=-(n-1)}^{n-1}\Big(1-\tfrac{\abs\tau}{n}\Big)\acvs_\tau,
\qquad
n\Var(\bar X)\xrightarrow[n\to\infty]{}\sum_{\tau=-\infty}^{\infty}\acvs_\tau .
\]
Hence $\Var(\bar X)\to0$ and $\bar X\toprob\mu$ (consistency).\\
\textit{Proof idea:} expand the double sum, collect by lag $\tau$, then apply
the \textbf{Ces\`aro summability theorem} to replace
$\sum(1-\tfrac{\abs\tau}{n})\acvs_\tau$ by $\sum\acvs_\tau$.\\
\textit{Why:} this is the prototype ``when can a sample average replace a
population average?''. Strong positive correlation inflates $\Var(\bar X)$
(e.g.\ AR(1) with $\phi\to1$): dependence, not just sample size, controls
estimation accuracy.
\end{proposition}

% ----------------------------------------------------------------------
% ----------------------------------------------------------------------
